\section{Principes de l'édition}

\subsection{La double page}

Chaque page du carnet occupe une double page : à droite le fac-similé, où les six pages écrites en travers sont pivotées d'un quart de tour et se lisent en tournant le livre ; à gauche la transcription, avec en marge un apparat aligné sur les lignes du manuscrit. Au bas, trois registres sans intitulé, séparés par des filets : remarques critiques, s'il y a lieu ; \textit{loci paralleli}, renvois (\textit{supra}, \textit{infra}) au passage précédent et au passage suivant où paraissent la même personne ou le même lieu (feuillet, ligne, page de l'édition), les noms les plus fréquents n'étant repris que dans l'index ; en retrait, commentaire historique. L'en-tête donne le mois et le feuillet restitué.

\subsection{La transcription}

La transcription est diplomatique : ligne à ligne, elle garde l'orthographe, les accents présents ou absents (\textit{Aout}, \textit{chateau}), les majuscules, la ponctuation et les tirets, et ne développe pas les abréviations. Elle rend la disposition : mois et montants rejetés à droite, soulignements, lettres suscrites, traits horizontaux, crochets et accolades groupant plusieurs lignes. Les lignes sont numérotées de cinq en cinq ; les traits et les blancs ne sont pas comptés. Signes employés :

\begin{center}
\small
\renewcommand{\arraystretch}{1.25}
\begin{tabular}{@{}>{\LJpolice}l p{92mm}@{}}
\toprule
\LJsl{25 Ma.} & mot souligné par Hippolyte Combes (\LJsll{double} : soulignement double)\\
M\LJsup{me}, 8\LJsup{h}, Lieu\LJsup{t} & lettres ou chiffres suscrits\\
\LJdel{venu} & mot biffé, resté lisible\\
\LJadd{très} & ajout dans l'interligne, à l'endroit où il s'insère\\
\LJinc{Sercus} & lecture incertaine\\
\LJill{} & mot illisible ; [\textit{3 mots illis.}] : plusieurs mots\\
\LJlac{} & lacune matérielle\\
\midrule
\multicolumn{2}{@{}l@{}}{\rule{25mm}{0.5pt}\hspace{3mm}trait horizontal tracé par Hippolyte Combes (séparation des journées, des comptes)}\\
\multicolumn{2}{@{}l@{}}{\tikz[baseline=1.4mm]{\draw[line width=0.5pt] (1.2mm,3.5mm) -- (0,3.5mm) -- (0,-1.5mm) -- (1.2mm,-1.5mm);}\hspace{2mm}crochet, et \tikz[baseline=0mm]{\draw[line width=0.45pt,decorate,decoration={brace,amplitude=1.1mm}] (0,3.5mm) -- (0,-1.5mm);}\hspace{1mm} accolade, groupant plusieurs lignes, imprimés à la hauteur des lignes groupées}\\
\bottomrule
\end{tabular}
\end{center}

Une surcharge (mot écrit sur un autre) est transcrite dans son état final ; la leçon recouverte est donnée dans l'apparat. Les mots anglais sont transcrits comme Hippolyte Combes les écrit. Pour que chaque ligne imprimée corresponde à une ligne du carnet, une ligne trop longue est réduite plutôt que coupée, et le corps et l'interligne varient d'une page à l'autre : plus petits pour une page dense ou chargée de notes, plus grands pour une page brève.

\subsection{L'apparat et les commentaires}\label{intro:apparat}

Les abréviations et les mots anglais courants ne sont ni développés ni traduits à chaque fois : la table des abréviations et le glossaire des termes anglais et militaires, en annexe, les réunissent. L'apparat développe les formes rares ou ambiguës, traduit les mots anglais moins courants, justifie les lectures difficiles ou incertaines (avec, au besoin, l'autre lecture possible) et donne les ratures, surcharges et ajouts que le balisage ne rend pas. Le lemme est suivi d'un crochet fermant ; les remarques sur une même ligne sont séparées par un trait vertical.

Le registre critique ne retient que ce que le fac-similé ne montre pas : l'ordre restitué des feuillets, un verso resté blanc (ainsi le f.~6~v\textsuperscript{o}), la continuité d'une journée d'une page à l'autre ou sa date. Il est le plus souvent vide.

Le commentaire historique, bref, identifie lieux, personnes et unités, et situe la page dans l'histoire de l'armée britannique sur la Somme et des interprètes militaires français. Il donne ses sources et signale les identifications incertaines. Il cite le carnet sous sa forme diplomatique ; dans l'introduction et les introductions de chapitre, les citations du carnet sont régularisées (abréviations développées, accents rétablis, fautes évidentes corrigées), la forme exacte se lisant au feuillet indiqué.

\subsection{La division en chapitres}

Le carnet ne comporte aucune division : c'est une suite continue de notes datées. Les sept chapitres de l'édition, découpage éditorial, suivent les déplacements d'Hippolyte Combes : de la Flandre à la Lys, de la Lys à la Somme, la première permission, l'Ancre à l'automne de 1916, l'hiver à Acheux, la seconde permission, Hérissart au printemps de 1917. Chaque chapitre s'ouvre sur une introduction et une carte.
